\chapter{Foundations: Mathematics}
\label{chap:foundations}

\section{Linear maps, directions and information}
\label{sec:foundations-linear-maps-directions-and-information}
\label{sec:la-vector-geometry-and-linear-independence}
\label{sec:la-matrices-as-linear-transformations}

For ML, the useful geometric question is usually not ``what is a vector?'' but ``which directions contain distinct information, and what does a transformation preserve or destroy?''
For nonzero vectors, the dot product separates magnitude from directional alignment:
\[
\underbrace{x^\top y}_{\text{magnitude + alignment}}
=
\|x\|\,\|y\|\underbrace{\cos\phi}_{\text{directional alignment}},
\qquad
x\perp y\iff x^\top y=0.
\]
Cosine similarity removes magnitude; the raw dot product does not.
This distinction matters in embeddings and similarity search when vector norm may or may not carry task information.
\label{subsec:la-dot-product}
\label{subsec:la-orthogonality}


\subsection{Independence, span, rank and null space}
\label{subsec:la-linear-independence}
\label{subsec:la-span-and-rank}
\label{subsec:la-image-space-and-null-space}
\label{subsec:la-singular-matrices}

\textbf{Linear independence} means a direction cannot be reconstructed from the others; the \textbf{span} is everything their combinations can represent; \textbf{rank} counts independent surviving directions.
For a linear map $y=Ax$, rank--nullity gives the information-capacity picture directly:
\begin{equation}
    \underbrace{N_{\text{input dimensions}}}_{\text{available directions}}
    =
    \underbrace{\operatorname{rank}(A)}_{\text{directions that survive}}
    +
    \underbrace{\dim\ker(A)}_{\text{directions destroyed}}
.
\end{equation}
The image/column space is what the map can output; the null space contains input perturbations that become invisible.
Redundant features and low-rank parameterisations are both statements about how many independent directions are actually needed.
This is the prerequisite for PCA in \cref{sec:representation-pca-an-eigenvector-view} and LoRA in \cref{sec:adaptation-lora-a-parameter-efficient-adaptation-mechanism}.

\begin{calloutmemoryhook}
Rank counts independent directions (columns of matrix); the null space counts directions the map deletes.
\end{calloutmemoryhook}

An inverse exists only when a square linear map loses no direction and covers its output space.
In numerical ML code, solve a linear system or use a factorisation rather than explicitly forming $A^{-1}$ unless the inverse itself is required.
\label{subsec:la-inverse-of-a-matrix}
\label{subsec:la-determinant}

\section{Eigenstructure and low-rank structure}
\label{sec:foundations-eigenstructure-and-low-rank-structure}
\label{sec:la-eigenstructure-and-matrix-properties}

\subsection{Eigenvectors, eigenvalues and PSD structure}
\label{subsec:la-eigenvectors-and-eigenvalues}
\label{subsec:la-positive-semi-definite-psd-matrices}

An eigenvector is a direction a linear map does not rotate away from itself:
\[
A\underbrace{v}_{\text{preserved direction}}
=
\underbrace{\lambda_{\text{along }v}}_{\text{scale on that direction}}v.
\]
The same mental model appears in covariance/PCA, recurrent dynamics and curvature: eigenvectors identify special directions; eigenvalues say how strongly the operator acts along them.
PCA uses covariance eigenstructure explicitly in \cref{sec:representation-pca-an-eigenvector-view}; the optimisation consequence reappears with Hessian curvature in \cref{sec:optimisation-newton-s-method-and-other-second-order-methods}.

A symmetric matrix is positive semidefinite (PSD) when
\[
\underbrace{x^\top A x}_{\text{quadratic response along }x}\ge 0\qquad\text{for every }x.
\]
For symmetric $A$, this is equivalent to nonnegative eigenvalues.
Covariance matrices are PSD.
\begin{calloutmemoryhook}
A PSD Hessian means convex space.
\end{calloutmemoryhook}

\subsection{SVD as the rectangular low-rank view}
\label{subsec:foundations-svd}

For a rectangular matrix, singular value decomposition exposes the same directional story without requiring a square matrix:
\[
A
=
\underbrace{U}_{\text{output directions}}
\underbrace{\Sigma}_{\text{direction-wise scales}}
\underbrace{V^\top}_{\text{input directions}}.
\]
Large singular values identify directions the map transmits strongly; truncating small ones gives a low-rank approximation.
That is the common structural idea behind compression, PCA implementations and parameter-efficient low-rank updates such as LoRA (\cref{sec:adaptation-lora-a-parameter-efficient-adaptation-mechanism}).

\section{Probability, uncertainty and estimation}
\label{sec:foundations-pmfs-and-pdfs}

For a discrete random variable, a PMF assigns probability mass and sums to one; for a continuous variable, a PDF integrates to one and probabilities are areas under the density.
Conditional and joint distributions connect through $p(x,y)=p(y\mid x)p(x)$; marginalisation removes variables by summing or integrating them out.

\subsection{Expectation, variance and covariance}
\label{sec:foundations-expectation-variance-and-covariance}

Expectation is the distribution-weighted average.
Variance and covariance capture spread and linear co-variation:
\begin{align}
\Var(X)
&=\E\!\left[(X-\E[X])^2\right],\\
\Cov(X,Y)
&=\E\!\left[(X-\E[X])(Y-\E[Y])\right],\\
\rho_{XY}
&=\frac{\Cov(X,Y)}{\sqrt{\Var(X)\Var(Y)}}.
\end{align}
Zero covariance is weaker than independence; independence implies zero covariance when the relevant moments exist.
For vector-valued data, the covariance matrix $\Sigma=\E[(x-\E[x])(x-\E[x])^\top]$ packages pairwise linear co-variation and feeds directly into PCA/whitening in \cref{sec:representation-pca-an-eigenvector-view,sec:representation-whitening-pca-and-zca}.


\subsection{Bayes, likelihood and priors}
\label{sec:foundations-bayes-theorem}

\begin{calloutmemorise}
\[
\boxed{P(A\mid B)=\frac{P(B\mid A)P(A)}{P(B)}}.
\]
\end{calloutmemorise}
Bayes' rule says a posterior combines prior belief with how compatible the observation is under each hypothesis.
For rare events, the prior/base rate can dominate the interpretation of a positive test; this is the probability-side version of precision in \cref{sec:evaluation-precision-recall-and-f1}.

\begin{calloutexample}
If prevalence is $1\%$, sensitivity is $95\%$, and the false-positive rate is $5\%$, then $P(\text{fault}\mid +)\approx16\%$: a sensitive detector can still produce mostly false alarms when the event is rare.
\end{calloutexample}

\label{sec:foundations-maximum-likelihood-and-maximum-a-posteriori-estimation}
For $N_{\text{examples}}$ conditionally independent observations, maximum likelihood chooses parameters that maximise summed log-likelihood:
\[
\hat\theta_{\mathrm{MLE}}
=
\argmax_{\theta}\sum_{i=1}^{N_{\text{examples}}}\log p(x_i\mid\theta).
\]
MAP adds a prior preference directly to the same objective:
\[
\hat\theta_{\mathrm{MAP}}
=
\argmax_{\theta}
\left[
\underbrace{\log p(D\mid\theta)}_{\text{fit to observations}}
+
\underbrace{\log p(\theta)}_{\text{parameter preference}}
\right].
\]
Gaussian-like priors produce L2-style penalties; Laplace-like priors produce L1-style penalties, linking Bayesian preference to regularisation in \cref{sec:stability-regularisation-l1-and-l2} and inductive bias in \cref{sec:nfl-bayesian-priors-a-direct-form-of-inductive-bias}.

\subsection{Entropy and distribution mismatch}
\label{sec:foundations-entropy}

\begin{calloutmemorise}
\[
\boxed{H(X)=-\sum_x p(x)\log p(x)}.
\]
\end{calloutmemorise}
Entropy measures uncertainty in one distribution.
Cross-entropy measures coding/surprise when observations follow $p$ but predictions use $q$; KL is the extra cost relative to using $p$ itself:
\label{sec:foundations-cross-entropy-kl-divergence-and-jensen-shannon-divergence}
\[
\underbrace{H(p,q)}_{\text{cross-entropy}}
=
\underbrace{H(p)}_{\text{irreducible uncertainty}}
+
\underbrace{D_{\mathrm{KL}}(p\|q)}_{\text{mismatch penalty}}.
\]
KL is asymmetric and can be infinite when $q$ assigns zero probability where $p$ does not.
Jensen--Shannon divergence symmetrises this comparison through the mixture $(p+q)/2$ and is bounded.
Cross-entropy becomes a classification objective in \cref{sec:objectives-classification-losses-and-logits}; entropy reduction also drives common tree-split criteria in \cref{subsec:trees-learning-a-split}.

\section{Local Derivatives and Curvature}
\label{sec:la-multivariable-calculus-and-optimisation}

\subsection{Jacobian: local input-to-output sensitivity}
\label{subsec:la-jacobian}

For a vector-valued map $f:x\mapsto y$, the local first-order response is
\[
\Delta y
\approx
\underbrace{\frac{\partial y}{\partial x}}_{\text{Jacobian: local sensitivity map}}\Delta x.
\]
The Jacobian tells which input directions are amplified, suppressed or locally invisible; its products are the algebra underneath backpropagation and recurrent gradient propagation.
See \cref{sec:optimisation-backpropagation-how-gradients-are-calculated,sec:stability-weight-initialisation-start-in-a-trainable-region} for those consequences rather than reteaching them here.

\begin{calloutmemoryhook}
The Jacobian maps a small input perturbation to its first-order output perturbation.
\end{calloutmemoryhook}

\subsection{Hessian: local curvature of a scalar objective}
\label{subsec:la-hessian}

For a scalar loss, the second derivative matrix controls the local quadratic correction:
\[
\mathcal L(\theta+\Delta)
\approx
\mathcal L(\theta)
+
\nabla_\theta\mathcal L(\theta)^\top\Delta
+
\frac12\Delta^\top
\underbrace{\nabla_\theta^2\mathcal L(\theta)}_{\text{Hessian: local curvature}}\Delta.
\]
Large positive Hessian eigenvalues mean sharp positive curvature; small magnitudes mean flat directions; negative eigenvalues identify locally downhill curvature directions around a saddle/non-minimum.
Newton and curvature-aware methods use this structure in \cref{sec:optimisation-newton-s-method-and-other-second-order-methods}.

\subsection{Convexity}
\label{sec:foundations-convexity-and-why-it-matters-for-optimisation}

A convex objective has no suboptimal local minima: averaging two parameter points cannot raise the function above the corresponding average of their objective values.
For twice-differentiable objectives on a convex domain, a PSD Hessian everywhere is the useful local test.
Deep-network losses are generally nonconvex, so convex guarantees do not transfer automatically.

\section{Confidence intervals and standard error}
\label{sec:foundations-confidence-intervals}

For $N_{\text{independent observations}}$ with sample standard deviation $s$, the mean estimator has standard error
\[
\operatorname{SE}(\hat\mu)=\frac{s}{\sqrt{N_{\text{independent observations}}}}.
\]
A confidence interval is a sampling-procedure statement, not a posterior probability that the fixed parameter lies in the realised interval.
Use a Student-$t$ critical value for normally distributed observations with unknown variance and small samples; use resampling or a model matched to the dependence structure when independence is not credible.
Evaluation-specific consequences are revisited in \cref{sec:evaluation-cross-validation-and-fold-choice,sec:evaluation-model-selection-and-data-leakage}.

